\chapter{RL, Optimal Control, and Control as Inference}

\section*{학습 목표}

이 장은 demonstration을 넘어 policy를 개선하려면 어떤 수학적 언어가 필요한지 다룬다. Chapter 7이 expert trajectory를 모방하는 문제를 다루었다면, Chapter 8은 reward, cost, constraint, uncertainty, exploration, action distribution을 다룬다. VLA 시대에도 이 개념들은 사라지지 않는다. 단지 language instruction과 visual context가 state와 objective 위에 추가되었을 뿐이다.

\begin{enumerate}
  \item language-conditioned RL 문제를 MDP/POMDP 관점에서 정식화한다.
  \item return, value, Q function, Bellman relation을 VLA policy와 연결한다.
  \item finite-horizon optimal control과 MPC가 VLA controller layer에서 갖는 의미를 설명한다.
  \item maximum entropy RL과 control as inference가 action-distribution VLA를 이해하는 데 주는 관점을 정리한다.
\end{enumerate}

\section{Why imitation is not enough}

Behavior cloning은 expert가 잘하는 구간에서는 강력하다. 그러나 expert data가 충분하지 않거나 task objective가 명시적인 cost/constraint를 포함하면 imitation만으로는 부족하다 \cite{sutton2018,haarnoja2018sac,kumar2020cql}. 예를 들어 ``컵을 식기건조대에 놓아라''라는 task에서 우리는 성공 여부, 충돌 회피, 힘 제한, 시간, smoothness, recovery 성공률을 함께 보고 싶다. 이런 요구는 supervised loss만으로 표현하기 어렵다.

정책이 실제 환경과 상호작용하면서 성능을 개선하려면 다음 질문이 필요하다.

\begin{itemize}
  \item 어떤 outcome이 좋은가?
  \item 어떤 constraint를 절대 위반하면 안 되는가?
  \item exploration은 얼마나 허용 가능한가?
  \item offline robot data만으로 policy improvement를 할 수 있는가?
  \item VLA가 낸 action distribution은 어떤 cost나 reward와 일관되는가?
\end{itemize}

\section{Language-conditioned MDP}

가장 단순하게는 language-conditioned robot control을 MDP로 쓸 수 있다.
\begin{equation}
  \mathcal{M}_{\task}
  =
  (\mathcal{S},\mathcal{A},P,r_{\task},\gamma).
  \label{eq:language-mdp}
\end{equation}
여기서 $\task$는 language instruction 또는 task embedding이고, $r_{\task}$는 그 instruction에 따라 달라지는 reward function이다. policy는
\begin{equation}
  \act_t \sim \pi_{\theta}(\act\mid \obs_t,\conf_t,\task,\bel_t)
  \label{eq:language-policy}
\end{equation}
로 표현할 수 있다.

실제 로봇은 완전한 state $\state_t$를 보지 못하므로 POMDP가 더 정확하다. 이때 policy는 observation history 또는 belief를 사용한다.
\begin{equation}
  \bel_t = \eta p(\obs_t\mid \state_t)\int p(\state_t\mid \state_{t-1},\act_{t-1})\bel_{t-1}d\state_{t-1}.
  \label{eq:rl-belief-update}
\end{equation}
VLA에서 memory, recurrent state, video history, tactile history는 모두 $\bel_t$를 근사하는 방법으로 볼 수 있다.

\section{Return, value, and Q function}

policy $\pi_{\theta}$의 discounted return은 다음과 같다.
\begin{equation}
  J(\theta)
  =
  \mathbb{E}_{\pi_{\theta}}
  \left[
  \sum_{t=0}^{T}
  \gamma^{t} r_{\task}(\state_t,\act_t)
  \right].
  \label{eq:rl-return}
\end{equation}
value function은 state 또는 belief에서 기대되는 return이다.
\begin{equation}
  V^{\pi}(\state)
  =
  \mathbb{E}_{\pi}
  \left[
  \sum_{k=0}^{\infty}
  \gamma^{k}r(\state_{t+k},\act_{t+k})
  \mid \state_t=\state
  \right].
  \label{eq:value-function}
\end{equation}
Q function은 특정 action을 먼저 선택한 뒤의 기대 return이다.
\begin{equation}
  Q^{\pi}(\state,\act)
  =
  r(\state,\act)
  +
  \gamma
  \mathbb{E}_{\state'\sim P}
  [V^{\pi}(\state')].
  \label{eq:q-function}
\end{equation}

Bellman optimality relation은
\begin{equation}
  Q^{\star}(\state,\act)
  =
  r(\state,\act)
  +
  \gamma
  \mathbb{E}_{\state'\sim P}
  \left[
  \max_{\act'}Q^{\star}(\state',\act')
  \right]
  \label{eq:bellman-optimality}
\end{equation}
로 쓸 수 있다.

VLA 시스템에서 $Q$나 value model은 action을 직접 내는 policy와 다른 역할을 할 수 있다. 예를 들어 LLM/VLM planner가 여러 skill 후보를 만들고, value/affordance model이 ``이 robot이 지금 할 수 있는가''를 평가할 수 있다. 이 구조는 VLA를 단독 policy가 아니라 proposal과 scoring의 결합으로 본다.

\section{Reward is not just success}

robot task reward는 단순 success indicator만으로 충분하지 않을 때가 많다.
\begin{equation}
  r_t
  =
  r^{task}_t
  -
  \lambda_c c^{collision}_t
  -
  \lambda_f c^{force}_t
  -
  \lambda_u \|\ctrlcmd_t\|^2
  -
  \lambda_j c^{jerk}_t.
  \label{eq:robot-reward}
\end{equation}
여기서 task success는 positive term이고, collision, force, control effort, jerk는 penalty term이다. Language-conditioned reward model은 instruction과 observation을 보고 progress 또는 success를 추정할 수 있다.
\begin{equation}
  \hat{r}_t
  =
  R_{\psi}(\obs_t,\task,\act_t,\obs_{t+1}).
  \label{eq:learned-reward}
\end{equation}

하지만 learned reward는 곧바로 안전한 objective가 아니다. reward model이 misspecified되면 policy는 실제 task success가 아니라 reward model의 빈틈을 최적화할 수 있다. Real VLA에서는 reward model, success detector, safety monitor를 분리해서 기록해야 한다.

\section{Finite-horizon optimal control}

Optimal control은 dynamics와 cost가 주어졌을 때 control sequence를 계산한다. finite horizon 문제는
\begin{equation}
\begin{aligned}
  \ctrlcmd_{0:T-1}^{\star}
  =
  \arg\min_{\ctrlcmd_{0:T-1}}
  &\quad
  \sum_{t=0}^{T-1}
  \ell(\state_t,\ctrlcmd_t;\task)
  +
  \ell_T(\state_T;\task)\\
  \text{s.t.}
  &\quad
  \state_{t+1}=f(\state_t,\ctrlcmd_t),\\
  &\quad
  \state_t\in\mathcal{X},\quad
  \ctrlcmd_t\in\mathcal{U}.
\end{aligned}
\label{eq:finite-horizon-oc}
\end{equation}
MPC는 이 문제를 매 시점 다시 푸는 receding-horizon control이다. VLA가 내는 target이나 subgoal은 cost의 parameter가 될 수 있다.
\begin{equation}
  \ell(\state_t,\ctrlcmd_t;\act^{VLA}_t)
  =
  \|\bm{x}_{ee,t}-\bm{x}^{ref}(\act^{VLA}_t)\|_{Q}^{2}
  +
  \|\ctrlcmd_t\|_{R}^{2}.
  \label{eq:vla-mpc-cost}
\end{equation}

이 관점에서 VLA는 low-level torque를 직접 예측하지 않아도 된다. VLA는 objective, reference, constraint, terminal condition을 제공하고, optimal control layer가 dynamics와 limit을 처리할 수 있다.

\section{Constrained RL and safety}

로봇에서는 reward maximization만으로 충분하지 않다. constraint violation은 높은 reward로 보상될 수 없는 경우가 많다. constrained MDP는 다음 문제로 쓸 수 있다.
\begin{equation}
\begin{aligned}
  \max_{\pi}
  &\quad
  \mathbb{E}_{\pi}\left[\sum_t \gamma^t r_t\right]\\
  \text{s.t.}
  &\quad
  \mathbb{E}_{\pi}\left[\sum_t \gamma^t c^i_t\right]\le d_i,
  \qquad i=1,\ldots,m.
\end{aligned}
\label{eq:constrained-rl}
\end{equation}
여기서 $c^i_t$는 collision, force excess, workspace violation, human-proximity risk 같은 cost이다.

Real VLA에서는 이 constraint가 여러 layer에 걸쳐 구현된다. semantic safety filter는 instruction을 검사하고, geometric planner는 forbidden region을 피하며, controller는 torque/force limit을 지키고, runtime monitor는 stop/recover를 결정한다. RL objective는 이 stack을 대체하지 않는다.

\section{Maximum entropy RL}

Maximum entropy RL은 return뿐 아니라 policy entropy를 함께 최적화한다.
\begin{equation}
  J_{\mathrm{MaxEnt}}(\pi)
  =
  \mathbb{E}_{\pi}
  \left[
  \sum_{t}
  r(\state_t,\act_t)
  +
  \alpha \mathcal{H}(\pi(\cdot\mid \state_t))
  \right].
  \label{eq:maxent-rl}
\end{equation}
entropy term은 여러 가능한 action을 유지하게 하며, multimodal behavior를 다루는 데 유용하다.

VLA의 action generation에서도 이 관점은 중요하다. 조작 task에서 하나의 observation에 대해 하나의 정답 action만 있는 경우는 드물다. 여러 grasp, 여러 approach direction, 여러 recovery path가 가능하다. 따라서 action distribution을 유지하는 policy가 평균 action을 내는 regression policy보다 더 적합할 수 있다.

\section{Control as inference}

Control as inference는 optimal control을 probabilistic inference 문제로 본다. trajectory $\tau$에 대해 optimality variable $\mathcal{O}$를 도입하고,
\begin{equation}
  p(\mathcal{O}_t=1\mid \state_t,\act_t)
  \propto
  \exp(r(\state_t,\act_t))
  \label{eq:optimality-likelihood}
\end{equation}
라고 두면, 좋은 trajectory posterior는
\begin{equation}
  p(\tau\mid \mathcal{O}_{0:T}=1)
  \propto
  p(\tau)
  \exp\left(\sum_{t=0}^{T}r(\state_t,\act_t)\right).
  \label{eq:trajectory-posterior}
\end{equation}
가 된다.

이 관점은 diffusion policy나 flow-based action model을 이해하는 데 도움이 된다. policy는 단일 action을 회귀하는 함수가 아니라, task와 observation 조건 아래에서 좋은 trajectory 또는 action chunk의 distribution을 샘플링하는 모델이 될 수 있다.
\begin{equation}
  \act_{t:t+H}
  \sim
  p_{\theta}(\act_{t:t+H}\mid \obs_{\le t},\task).
  \label{eq:action-chunk-distribution}
\end{equation}

\section{Offline RL and robot data}

Real robot online RL은 비용과 위험이 크다. 따라서 VLA 시대에는 offline robot dataset에서 policy를 개선하려는 압력이 크다. offline RL objective는 dataset distribution 안에서 policy improvement를 시도한다.
\begin{equation}
  \theta^\star
  =
  \arg\max_{\theta}
  \mathbb{E}_{(\state,\act)\sim\mathcal{D}}
  [Q_{\phi}(\state,\pi_{\theta}(\state))]
  -
  \beta D(\pi_{\theta}(\cdot\mid\state),\pi_{\mathcal{D}}(\cdot\mid\state)).
  \label{eq:offline-rl-regularized}
\end{equation}
두 번째 항은 policy가 dataset support 밖으로 나가는 것을 막는 regularization으로 볼 수 있다.

VLA fine-tuning에서도 같은 문제가 있다. 거대한 model이 dataset에 없는 action을 그럴듯하게 생성할 수 있지만, 그것이 robot에서 안전하다는 보장은 없다. 따라서 offline VLA improvement에는 conservative policy update, validation rollout, simulator/hardware-in-the-loop evaluation이 필요하다.

\section{Evaluation metadata}

RL 또는 optimal-control component가 들어간 VLA 실험은 다음 항목을 기록해야 한다.

\begin{itemize}
  \item reward source: hand-coded, learned reward, success detector, human preference
  \item constraint set: collision, force, velocity, workspace, semantic safety
  \item online interaction budget and safety intervention count
  \item offline dataset coverage and out-of-distribution action filter
  \item policy update method: BC, offline RL, on-policy RL, MPC-guided update
  \item action distribution type: Gaussian, token, mixture, diffusion, flow, chunk
  \item evaluation split: offline validation, simulator rollout, real robot rollout
\end{itemize}

\begin{casebox}{Case study: reward model과 safety monitor를 분리해야 하는 이유}
learned reward가 task progress를 잘 예측해도 collision, force, human proximity 같은 hard constraint를 항상 보장하지는 않는다. Real VLA에서는 reward maximization과 safety filtering을 같은 scalar objective로 합치기보다, reward/cost/constraint/monitor를 분리해 기록하는 편이 검증 가능하다.
\end{casebox}

\chapterwork
{Sutton and Barto, SAC, CQL을 읽고 reward, entropy, conservative update가 VLA action distribution에 어떤 의미를 갖는지 연결하라 \cite{sutton2018,haarnoja2018sac,kumar2020cql}.}
{offline VLA improvement에서 dataset support 밖의 action이 위험한 이유는 무엇인가?}
{language-conditioned task 하나에 대해 reward, hard constraint, intervention rule, offline validation rule을 분리한 objective schema를 작성하라.}

\section{요약}

RL과 optimal control은 VLA의 반대편에 있는 낡은 이론이 아니다. 이들은 VLA가 목표, 제약, 불확실성, action distribution을 어떻게 다루어야 하는지 설명하는 수학적 언어다. Behavior cloning은 expert action을 모방하지만, RL과 optimal control은 무엇이 좋은 행동인지, 어떤 constraint를 지켜야 하는지, 어떤 action distribution이 합리적인지 묻는다. Real VLA는 이 관점을 safety와 data engine 안에 넣어야 한다.

다음 장에서는 language가 붙기 전부터 존재했던 image-to-action 계보, 즉 visuomotor policy를 다룬다.
